\scanpage{311}
\mspage{10}
fournit \deleted{\unsure{un}}

\uline{Prop. 2} Avec les notations \inserted{et hyp.} de prop. 1, avec $A$ ou $B$
additives \deleted{karoubiennes}
\[
 \text{a)}\quad\underline{\operatorname{Hom}}_{\text{comm. à }\varinjlim}(Ss,A)
 \xrightarrow{\approx}A^{\cdot}
\]
\[
 \text{b)}\quad\underline{\operatorname{Hom}}_{\text{comm. à }\varprojlim}(Ss^o,B)
 \xrightarrow{\approx}B_{\cdot}.
\]
Nous allons \deleted{préciser}\inserted{expliciter} cette correspondance du point
de vue de l'\uline{algèbre} universelle. Prenons d'abord le cas de b), avec
$B=(Ab)$ pour simplifier. \deleted{\unsure{avec $K_{\cdot}\in B$}}
\inserted{Donc il faut} préciser la relation entre
\[
 \underset{\text{commute aux }\varprojlim}{F:Ss^o\longrightarrow Ab}
 \quad\text{et}\quad K_{\cdot}=(K_n,d_n)\in Ab_{\cdot}
\]
qui se déterminent mutuellement. On désigne comme de juste par $K_*$ l'objet
de $Ab_*$ défini par $K_{\cdot}$. On a donc
\[
 F(X_*)\simeq\operatorname{Hom}_{\widehat{\boldsymbol\Delta}}(X_*,K_*)
 \simeq\operatorname{Hom}_{\widehat{\boldsymbol\Delta}_{ab}}(\mathbb Z(X_*),K_*)
\]
où $\unsure{\mathbb Z(X_*)|\{\mathbb Z(X_*)\}_*}$ désigne le ``faisceau abélien
libre engendré par $X_*$'' en l'occurrence
\[
 \mathbb Z(X_*)=(n\longmapsto\mathbb Z(X_n)=\mathbb Z^{(X_n)})
\]
Ce groupe \inserted{ab} $\tfrac12$ simplicial sera aussi noté $\mathbb C_*(X_*)$,
et $\mathbb C_{\cdot}(X_*)$ le complexe de chaînes associé \inserted{par DP}:
\deleted{On aura donc}
\[
 \left\{\begin{array}{l}
 \mathbb C_*(X_*)=\mathbb Z(X_*)=(n\mathrel{\unsure{\vdash|\mapsto}}\mathbb Z(X_n)=\mathbb Z^{(X_n)}\\
 \mathbb C_{\cdot}(X_*)=ND_{\cdot}(\mathbb C_*(X_*))
 \end{array}\right.
\]
Donc par Dold-Puppe on aura
\[
 F(X_*)\simeq
 \operatorname{Hom}_{\substack{\widehat{\boldsymbol\Delta}_{ab}\\\Vert\\Ab_*}}
 (\mathbb C_*(X_*),K_*)
 \underset{\text{par D.P.}}{\simeq}
 \operatorname{Hom}_{Ab_{\cdot}}(\mathbb C_{\cdot}(X_*),K_{\cdot})
\]
Donc le foncteur $F_{K_{\cdot}}=F\quad Ss^o\longrightarrow Ab$
\[
 X_*\longmapsto F_{K_{\cdot}}(X_*)=F(X_*)
\]
associé à $K_{\cdot}\in Ab_{\cdot}$ est donné par
\[
 \boxed{F_{\unsure{K_{\cdot}}}(X_*)=operatorname{Hom}_{\cdot}(\mathbb C_{\cdot}(X_*),K_{\cdot})}
\]
On en tire deux conclusions:

$1^\circ)$ Les valeurs sur $X_*$ des foncteurs envisagés $F$
($Ss^o\mathbin{\text{---}}Ab$, commutant aux $\varprojlim$) ne dépendent que de
$\mathbb C_*(X_*)\in Ab_*$ \inserted{(``abélianisé de $X_*$'')}, ou ce qui revient
au même, de $\mathbb C_{\cdot}(X_*)\in Ab_{\cdot}$;
\deleted{\unsure{donc le}} complexe de chaînes $\mathbb C_{\cdot}(X_*)_{\unsure{*}}$
peut être considéré comme ``l'invariant additif universel'' de $X_*$.
Quand on ``oublie $X_*$ au profit de $\mathbb C_{\cdot}(X_*)^{\unsure{*}}=C_{\cdot}\in Ab_{\cdot}$''
les foncteurs envisagés (en tant que foncteurs en $C_{\cdot}$) sont encore
\unsure{exactt|exactement} les foncteurs représent.bles sur $Ab_{\cdot}$.

[\uline{NB} Quand le foncteur $F$ a une structure plus riche, p.ex. d'une
multiplication $F\otimes F\mathbin{\text{---}}F$ i.e. est à valeurs dans les
$\mathbb Z$-\uline{algèbres}, \deleted{et \unsure{F|si}} cette structure sur $F(X_*)$
n'est pas reconstructible par
\margin{312}
